\documentclass[12pt]{article}
\usepackage{enumitem}
\usepackage{geometry}
\usepackage{amsmath, amssymb}
\geometry{margin=1in}
\begin{document}
\begin{center}
History Assessment 1 \
Total Marks: 40 \
Answer all questions.
\end{center}

\section*{Multiple Choice (20 marks)}
\begin{enumerate}[label=\arabic*.]
\item Who was part of the elite class in the Olmec culture?
\begin{enumerate}[label=(\alph*)]
    \item those who lived in the particularly rich farmlands
    \item those who lived near resource-privileged areas
    \item those who lived near the river where trading took place
    \item those who lived near the volcanic rocks used in tool making
\end{enumerate}
\item The earliest-known use of bronze is found in:
\begin{enumerate}[label=(\alph*)]
    \item South America.
    \item Mesopotamia.
    \item the Indus Valley.
    \item China.
\end{enumerate}
\item After 400,000 years ago, hominid evolution is characterized by:
\begin{enumerate}[label=(\alph*)]
    \item bipedalism and relatively stable brain size.
    \item a slight decrease in average brain size.
    \item a rapid increase in brain size.
    \item slowly increased brain size and the first stone tools.
\end{enumerate}
\item The evolutionary histories of different kinds of organisms and their relationships to each other are illustrated in a:
\begin{enumerate}[label=(\alph*)]
    \item diastema.
    \item phylogeny.
    \item mosaic.
    \item genera,
\end{enumerate}
\item The Minoans inhabited the \_\_\_\_\_\_\_\_\_ island of Crete.
\begin{enumerate}[label=(\alph*)]
    \item Mediterranean
    \item Egyptian
    \item Mesopotamian
    \item Indonesian
\end{enumerate}
\item Which of the following are techniques of trace element analysis?
\begin{enumerate}[label=(\alph*)]
    \item proton magnetometer and electrical resistivity
    \item accelerator mass spectrometry and photosynthesis pathways
    \item carbon isotope analysis and experimental replication
    \item neutron activation analysis and X- ray fluorescence
\end{enumerate}
\item The death of Qin Dynasty emperor Ying Zheng was marked by what elaborate memorial?
\begin{enumerate}[label=(\alph*)]
    \item An enormous mortuary temple filled with gold and gems.
    \item An enormous pyramid constructed out of millions of huge limestone blocks.
    \item The collective suicide of 60 members of the royal family.
    \item An 8,000-member life-sized terra cotta army.
\end{enumerate}
\item When did Mohenjo-daro develop into a complex urban center?
\begin{enumerate}[label=(\alph*)]
    \item during the Early Shang Dynasty, from 1,000 to 400 B.P.
    \item during the Qin Dynasty, from 3,500 to 3,000 B.P.
    \item during the Mature Harappan period, from 4,500 to 4,000 B.P.
    \item during the Early Harappan period, from 8,500 to 6,000 B.P.
\end{enumerate}
\item Ancestral Puebloan culture in Chaco Canyon was comprised of:
\begin{enumerate}[label=(\alph*)]
    \item around 75 substantial settlements and more than 300 smaller communities.
    \item more than 500 towns and villages with monumental architecture.
    \item between 30 and 40 farmsteads and three large pueblos.
    \item one large city state with more than 300 outlying farmsteads.
\end{enumerate}
\item The idea that processes like weathering and erosion are responsible for the present condition of the earth is known as:
\begin{enumerate}[label=(\alph*)]
    \item uniformitarianism.
    \item natural selection.
    \item creationism.
    \item adaptation.
\end{enumerate}
\end{enumerate}

\section*{True / False (20 marks)}
\textit{Answer True or False and justify briefly.}
\begin{enumerate}[label=\arabic*.]
\item True or False: The correct answer to 'The Yang-shao culture gave way to the Lung-Shan sometime after:' is '5,000 B.P.'.
\item True or False: The correct answer to 'During which time frame is there evidence for anatomically modern Homo sapiens?' is '200,000 to present'.
\item True or False: The correct answer to 'What do we know about the people who built Stonehenge based on the analysis of animal teeth from a nearby village?' is 'They brought animals to the site from outside of the region.'.
\item True or False: The correct answer to 'What was the climate of Greater Australia 20,000 years ago?' is 'temperate'.
\item True or False: The correct answer to 'Which best describes Stonehenge?' is 'Chavin ritual monument'.
\item True or False: The correct answer to 'The possibility that Homo erectus possessed an aesthetic sense and may have created "art" is suggested by:' is 'the size of their brains, with a mean brain size slightly larger than modern humans.'.
\item True or False: The correct answer to 'The New World, "discovered" by Christopher Columbus, was filled with approximately how many people?' is 'tens of millions'.
\item True or False: The correct answer to 'Pompeii was unusually well preserved because:' is 'the layer of volcanic pumice created a seal over the remains.'.
\item True or False: The correct answer to 'The early Mesolithic Maglemosian culture was adapted to:' is 'an equatorial environment.'.
\item True or False: The correct answer to 'Opportunistic foraging involves:' is 'reliance on steady sources of food such as shellfish.'.
\end{enumerate}

\vfill
\noindent\textit{End of Paper}
\end{document}